% ==== III ====
\yoschapter{III}{The Orthogonal Projection and F. \textsc{Riesz}' Representation Theorem}
\label{chap:III}


\yossection{1}{The Orthogonal Projection}
\label{sec:III.1}

In a pre-Hilbert space, we can introduce the notion of orthogonality of two vectors. Thanks to this fact, a Hilbert space may be identified with its dual space, i.e., the space of bounded linear functionals. This result is the representation theorem of F. \textsc{Riesz} \cite{rickart272}, and the whole theory of Hilbert spaces is founded on this theorem. 

\begin{definition}\label{def:III.1.Definition.1} Let $x, y$ be vectors of a pre-Hilbert space $X$. We say that $x$ is orthogonal to $y$ and write $x \perp y$ if $(x, y) = 0$; if $x \perp y$ then $y \perp x$, and $x \perp x$ iff $x = 0$. Let $M$ be a subset of a pre-Hilbert space $X$. We denote by $M^{\perp}$ the totality of vectors $\in X$ orthogonal to every vector $m$ of $M$. 
\end{definition}

\begin{theorem}\label{thm:III.1.Theorem.1} Let $M$ be a closed linear subspace of a Hilbert space $X$. Then $M^{\perp}$ is also a closed linear subspace of $X$, and $M^{\perp}$ is called the orthogonal complement of $M$. Any vector $x \in X$ can be decomposed uniquely in the form 
\begin{equation}\label{eq:III.1.1}
x = m + n, \text{ where } m \in M \text{ and } n \in M ^{\perp}.
\end{equation}
\end{theorem}

The element $m$ in \eqref{eq:III.1.1} is called the orthogonal projection of $x$ upon $M$ and will be denoted by $P_M x; P_M$ is called the projection operator or the projector upon $M$. We have thus, remembering that $M \subseteq (M^\perp)^\perp$, 
\begin{equation}\tag{1'}\label{eq:III.1.1'}
x = P _{M} x + P _{M \perp} x, \text{ that is, } I = P _{M} + P _{M \perp}.
\end{equation}

\begin{proof}
 The linearity of $M^{\perp}$ is a consequence of the linearity in $x$ of the scalar product $(x, y)$. $M^{\perp}$ is closed by virtue of the continuity of the scalar product. The uniqueness of the decomposition \eqref{eq:III.1.1} is clear since a vector orthogonal to itself is the zero vector. 

To prove the possibility of the decomposition \eqref{eq:III.1.1}, we may assume that $M \neq X$ and $x \notin M$, for, if $x \in M$, we have the trivial decomposition with $m = x$ and $n = 0$. Thus, since $M$ is closed and $x \notin M$, we have 
\[
d = \inf\limits _{m \in M} \| x - m \| > 0.
\]

Let $\{m_n\} \subseteq M$ be a minimizing sequence, i.e., $\lim\limits_{n\to \infty}\|x - m_n\| = d$. Then $\{m_n\}$ is a Cauchy sequence. For, by $\|a + b\|^2 + \|a - b\|^2 = 2(\|a\|^2 + \|b\|^2)$ valid in any pre-Hilbert space (see \eqref{eq:I.5.1} in Chapter~\ref{chap:I}, \ref{sec:I.5}), we obtain 
\[
\begin{aligned}
&{\| m _{k} - m _{n} \| ^{2} = \| (x - m _{n}) - (x - m _{k}) \| ^{2} = 2 (\| x - m _{n} \| ^{2} + \| x - m _{k} \| ^{2})} \\
&{\qquad - \| 2 x - m _{n} - m _{k} \| ^{2}} {\qquad = 2 (\| x - m _{n} \| ^{2} + \| x - m _{k} \| ^{2}) - 4 \| x - (m _{n} + m _{k}) / 2 \| ^{2}} \\
&{\qquad \leqq 2 (\| x - m _{n} \| ^{2} + \| x - m _{k} \| ^{2} - 4 d ^{2})} {\qquad (\text{ since } (m _{n} + m _{k}) / 2 \in M)} \\
&{\qquad \rightarrow 2 (d ^{2} + d ^{2}) - 4 d ^{2} = 0 \text{ as } k, n \rightarrow \infty .}
\end{aligned}
\]

By the completeness of the Hilbert space $X$, there exists an element $m \in X$ such that $\slim\limits_{n \to \infty} m_n = m$. We have $m \in M$ since $M$ is closed. Also, by the continuity of the norm, we have $\|x - m\| = d$. 

Write $x = m + (x - m)$. Putting $n = x - m$, we have to show that $n \in M^{\perp}$. For any $m' \in M$ and any real number $\alpha$, we have $(m + \alpha m') \in M$ 
and so 
\[
\begin{aligned}
&d ^{2}  \leqq \| x - m - \alpha m ^{\prime} \| ^{2} = (n - \alpha m ^{\prime}, n - \alpha m ^{\prime}) = \| n \| ^{2} - \alpha (n, m ^{\prime}) - \alpha (m ^{\prime}, n) + \alpha^{2} \| m ^{\prime} \| ^{2}.
\end{aligned}
\]

This gives, since $\| n \| = d$, $0 \leqq -2\alpha \operatorname{Re}(n, m') + \alpha^2 \| m' \|^2$ for every real $\alpha$. Hence $\operatorname{Re}(n, m') = 0$ for every $m' \in M$. Replacing $m'$ by $im'$, we obtain $\operatorname{Im}(n, m') = 0$ and so $(n, m') = 0$ for every $m' \in M$. 
\end{proof}

\begin{corollaryu} For a closed linear subspace $M$ of a Hilbert space $X$, we have $M = M^{\perp \perp} = (M^{\perp})^{\perp}$. 
\end{corollaryu}

\begin{theorem}\label{thm:III.1.Theorem.2} The projector $P = P_{M}$ is a bounded linear operator such that 
\begin{align}
P& = P ^{2} \text{(idempotent property of } P), \label{eq:III.1.2} \\
(P x, y)& = (x, P y) \text{ (symmetric property of } P). \label{eq:III.1.3}
\end{align}
\end{theorem}

Conversely, a bounded linear operator $P$ on a Hilbert space $X$ into $X$ satisfying \eqref{eq:III.1.2} and \eqref{eq:III.1.3} is a projector upon $M = R(P)$. 

\begin{proof}
 \eqref{eq:III.1.2} is clear from the definition of the orthogonal projection. We have, by \eqref{eq:III.1.1} and $P_M x \perp P_M \perp y$, 
\[ \begin{aligned}
&(P _{M} x, y)  = (P _{M} x, P _{M} \bot y + P _{M} y) = (P _{M} x, P _{M} y) = (P _{M} x + P _{M} \bot x, P _{M} y) = (x, P _{M} y).
\end{aligned}
\]

Next let $y = x + z$, $x \in M$, $z \in M^{\perp}$ and $w = u + v$, $u \in M$, $v \in M^{\perp}$, then $y + w = (x + u) + (z + v)$ with $(x + u) \in M$, $(z + v) \in M^{\perp}$ and so, by the uniqueness of the decomposition \eqref{eq:III.1.1}, $P_M(y + w) = P_My + P_Mw$; similarly we obtain $P_M(xy) = \alpha P_My$. The boundedness of the operator $P_M$ is proved by 
\[
\begin{aligned} & {\| x \| ^{2} = \| P _{M} x + P _{M \perp} x \| ^{2} = (P _{M} x + P _{M \perp} x, P _{M} x + P _{M \perp} x)} \\ & {\quad = \| P _{M} x \| ^{2} + \| P _{M \perp} x \| ^{2} \geqq \| P _{M} x \| ^{2}.} \end{aligned}
\]

Thus, in particular, we have 
\begin{equation}\label{eq:III.1.4}
\left\| P _{M} \right\| \leqq 1.
\end{equation}

The converse part of the Theorem is proved as follows. The set $M = R(P)$ is a linear subspace, since $P$ is a linear operator. The condition $x \in M$ is equivalent to the existence of a certain $y \in X$ such that $x = P \cdot y$, and this in turn is equivalent, by \eqref{eq:III.1.2}, to $x = Py = P^2y = Px$. Therefore $x \in M$ is equivalent to $x = Px$. $M$ is a closed subspace; for, $x_n \in M$, $\slim x_n = y$ imply, by the continuity of $P$ and $x_n = Px_n$, $\slim x_n = \slim Px_n = Py$ so that $y = Py$. 

We have to show that $P = P_{M}$. If $x \in M$, we have $Px = x = P_{M} \cdot x$; and if $v \in M^{\perp}$, we have $P_{M}y = 0$. Moreover, in the latter case, $(Py, Py) = 6^{*}$ 
$(y, P^2 y) = (y, Py) = 0$ and so $Py = 0$. Therefore we obtain, for any $y \in X$, 
\[
\begin{aligned}
&P y  = P (P _{M} y + P _{M \perp} y) = P P _{M} y + P P _{M \perp} y = P _{M} y + 0, \text{ i.e., } P y = P _{M} y.
\end{aligned}
\]
\end{proof}

Another characterization of the projection operator is given by 

\begin{theorem}\label{thm:III.1.Theorem.3} $A$ bounded linear operator $P$ on a Hilbert space $X$ into $X$ is a projector iff $P$ satisfies $P = P^2$ and $\| P \| \leqq 1$. 
\end{theorem}

\begin{proof}
 We have only to prove the ``if'' part. Set $M = R(P)$ and $N = N(P) = \{y; P y = 0\}$. As in the proof of the preceding Theorem~\ref{thm:III.1.Theorem.2}, $M$ is a closed linear subspace and $x \in M$ is equivalent to $x = Px$. $N$ is also a closed linear subspace in virtue of the continuity of $P$. In the decomposition $x = Px + (I - P)x$, we have $Px \in M$ and $(I - P)x \in N$. The latter assertion is clear from $P(I - P) = P - P^{2} = 0$. 

We thus have to prove that $N = M^{\perp}$. For every $x \in X$, $y = Px - x \in N$ by $P^2 = P$. Hence, if, in particular, $x \in N^{\perp}$, then $Px = x + y$ with $(x, y) = 0$. It follows then that $\|x\|^2 \geqq \|Px\|^2 = \|x\|^2 + \|y\|^2$ so that $y = 0$. Thus we have proved that $x \in N^{\perp}$ implies $x = Px$, that is, $N^{\perp} \subseteq M = R(P)$. Let, conversely, $z \in M = R(P)$, so that $z = Pz$. Then we have the orthogonal decomposition $z = y + x$, $y \in N$, $x \in N^{\perp}$, and so $z = Pz = Py + Px = Px = x$, the last equality being already proved. This shows that $M = R(P) \subseteq N^{\perp}$. We have thus obtained $M = N^{\perp}$, and so, by $N = (N^{\perp})^{\perp}$, $N = M^{\perp}$. 
\end{proof}


\yossection{2}{``Nearly Orthogonal'' Elements}
\label{sec:III.2}

In general, we cannot define the notion of orthogonality in a normed linear space. However, we can prove the 

\begin{theoremn}[F. \textsc{Riesz} {\cite{rickart273}}] Let $X$ be a normed linear space, and $M$ be a closed linear subspace. Suppose $M \neq X$. Then there exists, for any $\varepsilon > 0$ with $0 < \varepsilon < 1$, an $x_{\varepsilon} \in X$ such that
\begin{equation}\label{eq:III.2.1}
\left\| x _{\varepsilon} \right\| = 1 \quad \text{ and } \quad \dist (x _{\varepsilon}, M) = \inf\limits _{m \in M} \left\| x _{\varepsilon} - m \right\| \geqq 1 - \varepsilon .
\end{equation}

The element $x_{\varepsilon}$ is thus ``nearly orthogonal'' to $M$. 

\end{theoremn}
\begin{proof}
 Let $y \in X - M$. Since $M$ is closed, $\dist(y, M) = \inf\limits_{m \in M} \| y - m \| = \alpha > 0$. Thus there exists an $m_{\varepsilon} \in M$ such that $\| y - m_{\varepsilon} \| \leqq \alpha \left(1 + \frac{\varepsilon}{1 - \varepsilon}\right)$. The vector $x_{\varepsilon} = (y - m_{\varepsilon}) / \| y - m_{\varepsilon} \|$ satisfies $\| x_{\varepsilon} \| = 1$ and 
\[ \begin{aligned}
&{\| x _{\varepsilon} - m \| = \| y - m _{\varepsilon} \| ^{-1} \| y - m _{\varepsilon} - \| y - m _{\varepsilon} \| \cdot m \| \geqq \| y - m _{\varepsilon} \| ^{-1} \alpha} {\quad \geqq \left(\frac {1}{1 - \varepsilon}\right) ^{-1} = 1 - \varepsilon .}
\end{aligned}
\]
\end{proof}

\begin{corollary}\label{cor:III.2.Corollary.1} Let there exist a sequence of closed linear subspaces $M_{n}$ of a normed linear space $X$ such that $M_{n} \subseteq M_{n+1}$ and $M_{n} \neq M_{n+1}$ ( $n = 1, 2, \ldots$). Then there exists a sequence $\{y_{n}\}$ such that 
\begin{equation}\label{eq:III.2.2}
y _{n} \in M _{n}, \| y _{n} \| = 1 \quad \text{ and } \quad \dist (y _{n+1}, M _{n}) \geqq 1 / 2 (n = 1, 2, \dots).
\end{equation}
\end{corollary}

\begin{corollary}\label{cor:III.2.Corollary.2} The unit sphere $S = \{x \in X; \| x \| \leqq 1\}$ of a $B$-space $X$ is compact iff $X$ is of finite dimension. 
\end{corollary}

\begin{proof}
 If $x_{1}, x_{2}, \ldots, x_{n}$ be a base for $X$, then the map $(\alpha_{1}, \alpha_{2}, \ldots, \alpha_{n}) \to \sum\limits_{j=1}^{n} \alpha_{j} x_{j}$ of $\mathbb{R}^{n}$ onto $X$ is surely continuous and so it is open by the open mapping theorem. This proves the ``if'' part. The ``only if'' part is proved as follows. Suppose $X$ is not of finite dimension. Then there exists, by the preceding Corollary~\ref{cor:III.2.Corollary.1}, a sequence $\{y_{n}\}$ satisfying the conditions: $\|y_{n}\|=1$ and $\|y_{m}-y_{n}\|\geqq1/2$ for $m>n$. This is a contradiction to the hypothesis that the unit sphere of $X$ is compact. 
\end{proof}


\yossection{3}{The Ascoli-Arzelà Theorem}
\label{sec:III.3}

To give an example of a relatively compact infinite subset of a $B$-space of infinite dimension, we shall prove the 

\begin{theoremn}[\textsc{Ascoli}-\textsc{Arzelà}] Let $S$ be a compact metric space, and $C(S)$ the $B$-space of (real- or) complex-valued continuous functions $x(s)$ normed by $\|x\| = \sup\limits_{s \in S} |x(s)|$. Then a sequence $\{x_n(s)\} \subseteq C(S)$ is relatively compact in $C(S)$ if the following two conditions are satisfied:
\begin{equation}\label{eq:III.3.1}
x _{n} (s) \text{ is equi-bounded (in } n), \text{ i.e., } \sup\limits _{n \geqq 1} \sup\limits _{s \in S} | x _{n} (s) | < \infty ,
\end{equation}
$x_{n}(s)$ is equi-continuous (in $n$), i.e., 
\begin{equation}\label{eq:III.3.2}
\lim\limits _{\delta \downarrow 0} \sup\limits _ {n \geqq 1, \dist (s ^{\prime}, s ^{\prime \prime}) \leqq \delta} | x _{n} (s ^{\prime}) - x _{n} (s ^{\prime \prime}) | = 0.
\end{equation}

\end{theoremn}
\begin{proof}
 A bounded sequence of complex numbers contains a convergent subsequence (the Bolzano-Weierstrass theorem). Hence, for fixed $s$, the sequence $\{x_{n}(s)\}$ contains a convergent subsequence. On the other hand, since the metric space $S$ is compact, there exists a countable dense subset $\{s_{n}\} \subseteq S$ such that, for every $\varepsilon > 0$, there exists a finite subset $\{s_{n}; 1 \leqq n \leqq k(\varepsilon)\}$ of $\{s_{n}\}$ satisfying the condition 
\begin{equation}\label{eq:III.3.3}
\sup\limits _{s \in S} \inf\limits _{1 \leqq j \leqq k (\varepsilon)} \dist (s, s _{j}) \leqq \varepsilon . \end{equation}

The proof of this fact is obtained as follows. Since $S$ is compact, it is totally bounded (see Chapter~\ref{chap:0}, \ref{sec:0.2}). Thus there exists, for any $\delta > 0$, a finite system of points $\in S$ such that any point of $S$ has a distance $\leqq \delta$ from some point of the system. Letting $\delta = 1, 2^{-1}, 3^{-1}, \ldots$ and collecting the corresponding finite systems of points, we obtain a sequence $\{s_n\}$ with the stated properties. 

We then apply the diagonal process of choice to the sequence $\{x_{n}(s)\}$, so that we obtain a subsequence $\{x_{n'}(s)\}$ of $\{x_{n}(s)\}$ which converges for $s = s_1, s_2, \ldots, s_k, \ldots$ simultaneously. By the equi-continuity of $\{x_{n}(s)\}$, there exists, for every $\varepsilon > 0$, a $\delta = \delta(\varepsilon) > 0$ such that $\text{ dis }(s', s'') \leqq \delta$ implies $|x_{n}(s') - x_{n}(s'')| \leqq \varepsilon$ for $n = 1, 2, \ldots$. Hence, for every $s \in S$, there exists a $j$ with $j \leqq k(\varepsilon)$ such that 
\begin{align}
\left| x _ {n ^{\prime}} (s) - x _ {m ^{\prime}} (s) \right|& \leqq \left| x _ {n ^{\prime}} (s) - x _ {n ^{\prime}} (s _{j}) \right| + \left| x _ {n ^{\prime}} (s _{j}) - x _ {m ^{\prime}} (s _{j}) \right| \notag \\
+ \left| x _ {m ^{\prime}} (s _{j}) - x _ {m ^{\prime}} (s) \right|& \leqq 2 \varepsilon + \left| x _ {n ^{\prime}} (s _{j}) - x _ {m ^{\prime}} (s _{j}) \right|. \notag
\end{align}

Thus $\lim\limits_{n,m\to \infty}\max\limits_s|x_{n'}(s) - x_{m'}(s)|\leqq 2\varepsilon$ and so $\lim\limits_{n,m\to \infty}\|x_{n'} - x_{m'}\| = 0.$ 
\end{proof}


\yossection{4}{The Orthogonal Base. Bessel's Inequality and Parseval's Relation}
\label{sec:III.4}

\begin{definition}\label{def:III.4.Definition.1} $A$ set $S$ of vectors in a pre-Hilbert space $X$ is called an orthogonal set, if $x \perp y$ for each pair of distinct vectors $x, y$ of $S$. If, in addition, $\|x\| = 1$ for each $x \in S$, then the set $S$ is called an orthonormal set. An orthonormal set $S$ of a Hilbert space $X$ is called a complete orthonormal system or an orthogonal base of $X$, if no orthonormal set of $X$ contains $S$ as a proper subset. 
\end{definition}

\begin{theorem}\label{thm:III.4.Theorem.1} $A$ Hilbert space $X$ (having a non-zero vector) has at least one complete orthonormal system. Moreover, if $S$ is any orthonormal set in $X$, there is a complete orthonormal system containing $S$ as a subset. 
\end{theorem}

Proof (by Zorn's Lemma). Let $S$ be an orthonormal set in $X$. Such a set surely exists; for instance, if $x \neq 0$, the set consisting only of $x / \|x\|$ is orthonormal. We consider the totality $\{S\}$ of orthonormal sets which contain $S$ as a subset. $\{S\}$ is partially ordered by writing $S_1 \prec S_2$ for the inclusion relation $S_1 \subseteq S_2$. Let $\{S'\}$ be a linearly ordered subsystem of $\{S\}$, then $\bigcup\limits_{S' \in \{S'\}} S'$ is an orthonormal set and an upper bound of $\{S'\}$. Thus, by Zorn's Lemma, there exists a maximal element $S_0$ of $\{S\}$. This orthonormal set $S_0$ contains $S$ and, by the maximality, it must be a complete orthonormal system. 

\begin{theorem}\label{thm:III.4.Theorem.2} Let $S = \{x_{\alpha}; \alpha \in A\}$ be a complete orthonormal system of a Hilbert space $X$. For any $f \in X$, we define its Fourier coefficients (with respect to $S$) 
\begin{equation}\label{eq:III.4.1}
f _{\alpha} = (f, x _{\alpha}).
\end{equation}
\end{theorem}

Then we have Parseval's relation 
\begin{equation}\label{eq:III.4.2}
\left\| f \right\| ^{2} = \sum\limits_{\alpha \in A} \left| f _{\alpha} \right| ^{2}.
\end{equation}

\begin{proof}
 We shall first prove Bessel's inequality 
\begin{equation}\tag{2'}\label{eq:III.4.2'}
\sum\limits_{\alpha \in A} | f _{\alpha} | ^{2} \leqq \| f \| ^{2}. \end{equation}

Let $\alpha_{1},\alpha_{2},\ldots ,\alpha_{n}$ be any finite system of $\alpha$'s. For any finite system of complex numbers $c_{\alpha_1},c_{\alpha_2},\ldots ,c_{\alpha_n}$, we have, by the orthonormality of $\{x_{\alpha}\}$, 
\begin{equation}\label{eq:III.4.3}
\begin{aligned}
&\left\| f - \sum\limits_{j=1} ^{n} c _ {\alpha_{j}} x _ {\alpha_{j}} \right\| ^{2}  = \left(f - \sum\limits_{j=1} ^{n} c _ {\alpha_{j}} x _ {\alpha_{j}}, f - \sum\limits_{j=1} ^{n} c _ {\alpha_{j}} x _ {\alpha_{j}}\right) \\
&= \| f \| ^{2} - \sum\limits_{j=1} ^{n} c _ {\alpha_{j}} \bar {f} _ {\alpha_{j}} - \sum\limits_{j=1} ^{n} \bar {c} _ {\alpha_{j}} f _ {\alpha_{j}} + \sum\limits_{j=1} ^{n} | c _ {\alpha_{j}} | ^{2} = \| f \| ^{2} - \sum\limits_{j=1} ^{n} | f _ {\alpha_{j}} | ^{2} + \sum\limits_{j=1} ^{n} | f _ {\alpha_{j}} - c _ {\alpha_{j}} | ^{2}.
\end{aligned}
\end{equation}

Hence the minimum of $\left\| f - \sum\limits_{j=1}^{n} c_{\alpha_j} x_{\alpha_j}\right\|^2$, for fixed $\alpha_1, \alpha_2, \ldots, \alpha_n$, is attained when $c_{\alpha_j} = f_{\alpha_j}$ ( $j = 1, 2, \ldots, n$). We have thus 
\begin{equation}\label{eq:III.4.4}
\left\| f - \sum\limits_{j=1} ^{n} f _ {\alpha_{j}} x _ {\alpha_{j}} \right\| ^{2} = \| f \| ^{2} - \sum\limits_{j=1} ^{n} | f _ {\alpha_{j}} | ^{2}, \text{ and hence } \sum\limits_{j=1} ^{n} | f _ {\alpha_{j}} | ^{2} \leqq \| f \| ^{2}.
\end{equation}

By the arbitrariness of $\alpha_{1},\alpha_{2},\ldots ,\alpha_{n}$, we see that Bessel's inequality \eqref{eq:III.4.2} is true, and $f_{\alpha}\neq 0$ for at most a countable number of $\alpha$'s, say $\alpha_{1},\alpha_{2},\ldots ,\alpha_{n},\ldots$ We then prove that $f = \slim\limits_{n\to \infty}\sum\limits_{j=1}^{n}f_{\alpha_j}x_{\alpha_j}$. First, the sequence $\left\{\sum\limits_{j=1}^{n}f_{\alpha_j}x_{\alpha_j}\right\}$ is a Cauchy sequence, since, by the orthonormality of $\{x_{\alpha}\}$, 
\[
\left| \sum\limits_{j=k} ^{n} f _ {\alpha_{j}} x _ {\alpha_{j}} \right| ^{2} = \left(\sum\limits_{j=k} ^{n} f _ {\alpha_{j}} x _ {\alpha_{j}}, \sum\limits_{j=k} ^{n} f _ {\alpha_{j}} x _ {\alpha_{j}}\right) = \sum\limits_{j=k} ^{n} | f _ {\alpha_{j}} | ^{2}
\]
which tends, by \eqref{eq:III.4.4} proved above, to $0$ as $k \to \infty$. We set $\slim\limits_{n \to \infty} \sum\limits_{j=1}^{n} f_{\alpha_j} x_{\alpha_j} = f'$, and shall prove that $(f - f')$ is orthogonal to every vector of $S$. By the continuity of the scalar product, we have 
\[
(f - f ^{\prime}, x _ {\alpha_{j}}) = \lim\limits _{n \rightarrow \infty} \left(f - \sum\limits_{k=1} ^{n} f _ {\alpha_{k}} x _ {\alpha_{k}}, x _ {\alpha_{j}}\right) = f _ {\alpha_{j}} - f _ {\alpha_{j}} = 0,
\]
and, when $\alpha \neq \alpha_{j}$ ( $j = 1, 2, \ldots$), 
\[
(f - f ^{\prime}, x _{\alpha}) = \lim\limits _{n \rightarrow \infty} \left(f - \sum\limits_{k=1} ^{n} f _{\alpha k} x _{\alpha k}, x _{\alpha}\right) = 0 - 0 = 0.
\]

Thus, by the completeness of the orthonormal system $S = \{x_{\alpha}\}$, we must have $(f - f') = 0$. Hence, by \eqref{eq:III.4.4} and the continuity of the norm, we have 
\[
0 = \lim\limits _{n \rightarrow \infty} \left\| f - \sum\limits_{j=1} ^{n} f _ {\alpha_{j}} x _ {\alpha_{j}} \right\| ^{2} = \| f \| ^{2} - \lim\limits _{n \rightarrow \infty} \sum\limits_{j=1} ^{n} | f _ {\alpha_{j}} | ^{2} = \| f \| ^{2} - \sum\limits_{\alpha \in A} | f _{\alpha} | ^{2}.
\]
\end{proof}

\begin{corollary}\label{cor:III.4.Corollary.1} We have the Fourier expansion 
\begin{equation}\label{eq:III.4.5}
f = \sum\limits_{j=1} ^{\infty} f _ {\alpha_{j}} x _ {\alpha_{j}} = \slim _{n \rightarrow \infty} \sum\limits_{j=1} ^{n} f _ {\alpha_{j}} x _ {\alpha_{j}}.
\end{equation}
\end{corollary}

\begin{corollary}\label{cor:III.4.Corollary.2} Let $l^2(A)$ be the space $L^2(A, \mathfrak{B}, m)$ where $m(\{\alpha\}) = 1$ for every point $\alpha$ of $A$. Then the Hilbert space $X$ is isometrically isomorphic to the Hilbert space $l^2(A)$ by the correspondence 
\begin{equation}\label{eq:III.4.6}
X \ni f \leftrightarrow \{f _{\alpha} \} \in l ^{2} (A)
\end{equation}
in the sense that 
\begin{equation}\label{eq:III.4.7}
(f + g) \leftrightarrow \left\{f _{\alpha} + g _{\alpha} \right\}, \beta f \leftrightarrow \left\{\beta f _{\alpha} \right\} \text{ and } \| f \| ^{2} = \| \left\{f _{\alpha} \right\} \| ^{2} = \sum\limits_{\alpha \in A} | f _{\alpha} | ^{2}.
\end{equation}
\end{corollary}

\begin{examplen} $\left\{\frac{1}{\sqrt{2\pi}} e^{int}; n = 0, \pm 1, \pm 2, \ldots\right\}$ is a complete orthonormal system in the Hilbert space $L^2(0, 2\pi)$.

\end{examplen}
\begin{proof}
 We have only to prove the completeness of this system. We have, by \eqref{eq:III.4.3}, 
\[
\begin{aligned}
\left\| f - \sum\limits_{j=-n} ^{n} \frac {1}{\sqrt {2 \pi}} c _{j} e ^{ijt} \right\| ^{2} \geqq \left\| f - \sum\limits_{j=-n} ^{n} \frac {1}{\sqrt {2 \pi}} f _{j} e ^{ijt} \right\| ^{2} &= \| f \| ^{2} - \sum\limits_{j=-n} ^{n} \frac {1}{2 \pi} | f _{j} | ^{2},
\\
\text{ where } f _{j} &= (f, e ^{ijt}). \end{aligned}
\]

If $f \in L^2(0, 2\pi)$ is continuous and with period $2\pi$, then the left hand side of the inequality above may be taken arbitrarily small by virtue of Weierstrass' trigonometric approximation theorem (see Chapter~\ref{chap:0}, \ref{sec:0.2}). Thus the set of all the linear combinations $\sum\limits_j c_j e^{ijt}$ is dense, in the sense of the norm, in the subspace of $L^2(0, 2\pi)$ consisting of all continuous functions with period $2\pi$. Such a subspace is also dense, in the sense of the norm, in the space $L^2(0, 2\pi)$. Therefore, any function $f \in L^2(0, 2\pi)$, orthogonal to all the functions of $\left\{\frac{1}{\sqrt{2\pi}} e^{int}\right\}$ must be a zero vector of $L^2(0, 2\pi)$. This proves that our system of functions $\frac{1}{\sqrt{2\pi}} e^{int}$ is a complete orthonormal system of $L^2(0, 2\pi)$. 
\end{proof}


\yossection{5}{E. \textsc{Schmidt}'s Orthogonalization}
\label{sec:III.5}

\begin{theoremn}[E. \textsc{Schmidt}'s orthogonalization] Given a finite or countably infinite sequence $\{x_{j}\}$ of linearly independent vectors of a pre-Hilbert space $X$. Then we can construct an orthonormal set having the same cardinal number as the set $\{x_{j}\}$ and spanning the same linear subspace as $\{x_{j}\}$.

\end{theoremn}
\begin{proof}
 Certainly $x_{1} \neq 0$. We define $y_{1}, y_{2}, \ldots$ and $u_{1}, u_{2}, \ldots$ recurrently as follows: 
\[
\begin{aligned}
&y_{1} = x_{1},  u_{1} = y_{1}/\|y_{1}\|, y_{2} = x_{2} - (x_{2}, u_{1})u_{1},  u_{2} = y_{2}/\|y_{2}\|, \cdots   \ddots \\
&y_{n+1} = x_{n+1} - \sum\limits_{j=1}^{n}(x_{n+1}, u_{j})u_{j},  u_{n+1} = y_{n+1}/\|y_{n+1}\|, \\
&\cdots   \ddots
\end{aligned}
\]

This process terminates if $\{x_j\}$ is a finite set. Otherwise it continues indefinitely. Observe that $y_n \neq 0$, because $x_1, x_2, \ldots, x_n$ are linearly independent. Thus $u_n$ is well defined. It is clear, by induction, that each $u_n$ is a linear combination of $x_1, x_2, \ldots, x_n$ and that each $x_n$ is a linear combination of $u_1, u_2, \ldots, u_n$. Thus the closed linear subspace spanned by the $u$'s is the same as that spanned by the $x$'s. 


We see, by $\|u_{1}\|=1$, that $y_{2}\perp u_{1}$ and hence $u_{2}\perp u_{1}$. Thus, by $\|u_{1}\|=1$, $y_{3}\perp u_{1}$ and hence $u_{3}\perp u_{1}$. Repeating the argument, we see that $u_{1}$ is orthogonal to $u_{2}, u_{3}, \ldots, u_{n}, \ldots$. Thus, by $\|u_{2}\|=1$, we have $y_{3}\perp u_{2}$ and so $u_{3}\perp u_{2}$. Repeating the argument, we finally see that $u_{k}\perp u_{m}$ whenever $k>m$. Therefore $\{u_{j}\}$ constitutes an orthonormal set.
\end{proof}

\begin{corollaryu} Let a Hilbert space $X$ be separable, i.e., let $X$ have a dense subset which is at most countable. Then $X$ has a complete orthonormal system consisting of an at most countable number of elements. 
\end{corollaryu}

\begin{proof}
 Suppose that an at most countable sequence $\{a_j\}$ of vectors $\in X$ be dense in $X$. Let $x_1$ be the first non-zero element in the sequence $\{a_j\}$, $x_2$ the first $a_i$ which is not in the closed subspace spanned by $x_1$, and $x_n$ the first $a_i$ which is not in the closed subspace spanned by $x_1, x_2, \ldots, x_{n-1}$. It is clear that the $a$'s and the $x$'s span the same closed linear subspace. In fact, this closed linear subspace is the whole space $X$, because the set $\{a_j\}$ is dense in $X$. Applying Schmidt's orthogonalization to $\{x_j\}$, we obtain an orthonormal system $\{u_j\}$ which is countable and spans the whole space $X$. 


This system $\{u_{j}\}$ is complete, since otherwise there would exist a non-zero vector orthogonal to every $u_{j}$ and hence orthogonal to the space $X$ spanned by $u_{j}$'s.
\end{proof}

\begin{examplen} Let $S$ be the interval $(a, b)$, and consider the real Hilbert space $L^{2}(S, \mathfrak{B}, m)$, where $B$ is the set of all Baire subsets in $(a, b)$. If we orthogonalize the set of monomials 
\[
1, s, s ^{2}, s ^{3}, \dots , s ^{n}, \dots ,
\]
we get the so-called Tchebyschev system of polynomials 
\[
P _{0} (s) = \text{ constant }, P _{1} (s), P _{2} (s), P _{3} (s), \dots , P _{n} (s), \dots
\]
which satisfies 
\[
\int_{a} ^{b} P _{i} (s) P _{j} (s) m (\dd s) = \delta_{ij} (= 0 \text{ or } 1 \text{ according as } i \neq j \text{ or } i = j).
\]

In the particular case when $a = -1$, $b = 1$ and $m(\dd s) = \dd s$, we obtain the Legendre polynomials; when $a = -\infty$, $b = \infty$ and $m(\dd s) = e^{-s^2} \dd s$, we obtain the Hermite polynomials and finally when $a = 0$, $b = \infty$ and $m(\dd s) = e^{-s} \dd s$, we obtain the Laguerre polynomials. 
\end{examplen}

It is easy to see that, when $-\infty < a < b < \infty$, the orthonormal system $\{P_{j}(s)\}$ is complete. For, we may follow the proof of the completeness of the trigonometric system (see the Example in the preceding section $4$); we shall appeal to Weierstrass' polynomial approximation theorem, in place of Weierstrass' trigonometric approximation theorem. As to the completeness proof of the Hermite or Laguerre polynomials, we refer the reader to G. \textsc{Szegő} \cite{szego313} or to K. \textsc{Yosida} \cite{yosida344}. 


\yossection{6}{F. \textsc{Riesz}' Representation Theorem}
\label{sec:III.6}

\begin{theoremn}[F. \textsc{Riesz}' representation theorem] Let $X$ be a Hilbert space and $f$ a bounded linear functional on $X$. Then there exists a uniquely determined vector $y_f$ of $X$ such that
\begin{equation}\label{eq:III.6.1}
f (x) = \left(x, y _{f}\right) \text{ for all } x \in X, \text{ and } \| f \| = \| y _{f} \|.
\end{equation}

Conversely, any vector $y \in X$ defines a bounded linear functional $f_y$ on $X$ by 
\begin{equation}\label{eq:III.6.2}
f _{y} (x) = (x, y) \text{ for all } x \in X, \text{ and } \| f _{y} \| = \| y \|.
\end{equation}

\end{theoremn}
\begin{proof}
 The uniqueness of $y_{f}$ is clear, since $(x,z)=0$ for all $x\in X$ implies $z=0$. To prove its existence, consider the null space $N=N(f)=\{x\in X;f(x)=0\}$ of $f$. Since $f$ is continuous and linear, $N$ is a closed linear subspace. The Theorem is trivial in the case when $N=X$; we take in this case, $y_{f}=0$. Suppose $N\neq X$. Then there exists a $y_{0}\neq0$ which belongs to $N^{\perp}$ (see Theorem~\ref{thm:III.1.Theorem.1} in Chapter~\ref{chap:III}, \ref{sec:III.1}). Define 
\begin{equation}\label{eq:III.6.3}
y _{f} = (\overline {{f (y _{0})}} / \| y _{0} \| ^{2}) y _{0}. \end{equation}

We will show that this $y_{f}$ meets the condition of the Theorem. First, if $x \in N$, then $f(x) = (x, y_{f})$ since both sides vanish. Next, if $x$ is of the form $x = \alpha y_{0}$, then we have 
\[
(x, y _{f}) = (\alpha y _{0}, y _{f}) = \left(\alpha y _{0}, \frac {\overline {{f (y _{0})}}}{\| y _{0} \| ^{2}} y _{0}\right) = \alpha f (y _{0}) = f (\alpha y _{0}) = f (x).
\]

Since $f(x)$ and $(x, y_{f})$ are both linear in $x$, the equality $f(x) = (x, y_{f})$, $x \in X$, is proved if we have proved that $X$ is spanned by $N$ and $y_{0}$. To show the last assertion, we write, remembering that $f(y_{f}) \neq 0$, 
\[
x = \left(x - \frac {f (x)}{f (y _{f})} y _{f}\right) + \frac {f (x)}{f (y _{f})} y _{f}.
\]

The first term on the right is an element of $N$, since 
\[
f \left(x - \frac {f (x)}{f \left(y _{f}\right)} y _{f}\right) = f (x) - \frac {f (x)}{f \left(y _{f}\right)} f \left(y _{f}\right) = 0.
\]

We have thus proved the representation $f(x) = (x, y_f)$. 

Therefore, we have 
\[
\left\| f \right\| = \sup\limits _{\left\| z \right\| \leqq 1} \left| f (x) \right| = \sup\limits _{\left\| z \right\| \leqq 1} \left| (x, y _{f}) \right| \leqq \sup\limits _{\left\| z \right\| \leqq 1} \left\| x \right\| \cdot \left\| y _{f} \right\| = \left\| y _{f} \right\|,
\]
and also $\|f\| = \sup\limits_{\|x\|\leqq 1}|f(x)| \geqq |f(y_f / \|y_f\|)| = \left(\frac{y_f}{\|y_f\|}, y_f\right) = \|y_f\|$. 

Hence we have proved the equality $\|f\| = \|y_f\|$. 

Finally, the converse part of the Theorem is clear from $|f_y(x)| = |(x, y)| \leqq \|x\| \cdot \|y\|$. 
\end{proof}

\begin{corollary}\label{cor:III.6.Corollary.1} Let $X$ be a Hilbert space. Then the totality $X'$ of bounded linear functionals on $X$ constitutes also a Hilbert space, and there is a norm-preserving, one-to-one correspondence $f \leftrightarrow y_{f}$ between $X'$ and $X$. By this correspondence, $X'$ may be identified with $X$ as an abstract set; but it is not allowed to identify, by this correspondence, $X'$ with $X$ as linear spaces, since the correspondence $f \leftrightarrow y_{f}$ is conjugate linear: 
\begin{equation}\label{eq:III.6.4}
\left(\alpha_{1} f _{1} + \alpha_{2} f _{2}\right) \leftrightarrow \left(\bar {\alpha} _{1} y _ {f _{1}} + \bar {\alpha} _{2} y _ {f _{2}}\right),
\end{equation}
where $\alpha_{1},\alpha_{2}$ are complex numbers. 
\end{corollary}

\begin{proof}
 It is easily verified that $X'$ is made into a Hilbert space by defining its scalar product through $(f_1, f_2) = (\overline{y_{f_1}, y_{f_2}})$, so that the statement of Corollary~\ref{cor:III.6.Corollary.1} is clear. 
\end{proof}

\begin{corollary}\label{cor:III.6.Corollary.2} Any continuous linear functional $T$ on the Hilbert space $X'$ is thus identified with a uniquely determined element $t$ of $X$ as follows: 
\begin{equation}\label{eq:III.6.5}
T (f) = f (t) \quad \text{ for all } \quad f \in X ^{\prime}.
\end{equation}
\end{corollary}

\begin{proof}
 Clear from the fact that the product of two conjugate linear transformations is a linear transformation. 
\end{proof}

\begin{definitionu} The space $X'$ is called the dual space of $X$. We can thus identify a Hilbert space $X$ with its second dual $X'' = (X')'$ in the above sense. This fact will be referred to as the reflexivity of Hilbert spaces. 
\end{definitionu}

\begin{corollary}\label{cor:III.6.Corollary.3} Let $X$ be a Hilbert space, and $X'$ its dual Hilbert space. Then, for any subset $F$ of $X'$ which is dense in the Hilbert space $X'$, we have 
\begin{equation}\label{eq:III.6.6}
\left\| x _{0} \right\| = \sup\limits _{f \in F, \| f \| \leqq 1} \left| f (x _{0}) \right|, x _{0} \in X.
\end{equation}
\end{corollary}

\begin{proof}
 We may assume that $x_0 \neq 0$, otherwise the formula \eqref{eq:III.6.6} is trivial. We have $(x_0, x_0 / \|x_0\|) = \|x_0\|$, and so there exists a bounded linear functional $f_0$ on $X$ such that $\|f_0\| = 1$, $f_0(x_0) = \|x_0\|$. Since $f(x_0) = (x_0, y_f)$ is continuous in $y_f$, and since the correspondence $f \leftrightarrow y_f$ is norm-preserving we see that \eqref{eq:III.6.6} is true by virtue of the denseness of $F$ in $X'$. 
\end{proof}

\begin{remarku} Hilbert's original definition of the ``Hilbert space'' is the space $(l^2)$. See his paper \cite{hilbert96}. It was J. von \textsc{Neumann} \cite{neumann228} who gave an axiomatic definition (see Chapter~\ref{chap:I}, \ref{sec:I.9}) of the Hilbert space assuming that the space is separable. F. \textsc{Riesz} \cite{rickart272} proved the above representation theorem without assuming the separability of the Hilbert space. In this paper, he stressed that the whole theory of the Hilbert space may be founded upon this representation theorem. 
\end{remarku}


\yossection{7}{The Lax-Milgram Theorem}
\label{sec:III.7}

Of recent years, it has been proved that a variant of F. \textsc{Riesz}' representation theorem, formulated by P. \textsc{Lax} and A. N. \textsc{Milgram} \cite{lax171}, is a useful tool for the discussion of the existence of solutions of linear partial differential equations of elliptic type. 

\begin{theoremn}[\textsc{Lax}-\textsc{Milgram}] Let $X$ be a Hilbert space. Let $B(x, y)$ be a complex-valued functional defined on the product Hilbert space $X \times X$ which satisfies the conditions:

Sesqui-linearity, i.e., 
\begin{equation}\label{eq:III.7.1}
\begin{aligned}
B \left(\alpha_{1} x _{1} + \alpha_{2} x _{2}, y\right) &= \alpha_{1} B \left(x _{1}, y\right) + \alpha_{2} B \left(x _{2}, y\right) \text{ and }
\\
B \left(x, \beta_{1} y _{1} + \beta_{2} y _{2}\right) &= \bar {\beta} _{1} B \left(x, y _{1}\right) + \bar {\beta} _{2} \dot {B} \left(x, y _{2}\right),
\end{aligned}
\end{equation}

Boundedness, i.e., there exists a positive constant $\gamma$ such that 
\begin{equation}\label{eq:III.7.2}
\left| B (x, y) \right| \leqq \gamma \left\| x \right\| \cdot \left\| y \right\|,
\end{equation}

Positivity, i.e., there exists a positive constant $\delta$ such that 
\begin{equation}\label{eq:III.7.3}
B (x, x) \geqq \delta \| x \| ^{2}.
\end{equation}

Then there exists a uniquely determined bounded linear operator $S$ with a bounded linear inverse $S^{-1}$ such that 
\begin{equation}\label{eq:III.7.4}
\begin{minipage}[b]{0.92\linewidth}
$(x,y) = B(x,Sy)$ whenever $x$ and $y\in X$, and $\| S\| \leqq \delta^{-1},\| S^{-1}\| \leqq \gamma$.
\end{minipage}
\end{equation}


\end{theoremn}
\begin{proof}
 Let $D$ be the totality of elements $y \in X$ for which there exists an element $y^*$ such that $(x, y) = B(x, y^*)$ for all $x \in X$. $D$ is not void, because $0 \in D$ with $0^* = 0 \cdot y^*$ is uniquely determined by $y$. For, if $w$ be such that $B(x, w) = 0$ for all $x$, then $w = 0$ by $0 = B(w, w) \geqq \delta \|w\|^2$. By the sesqui-linearity of $(x, y)$ and $B(x, y)$, we obtain a linear operator $S$ with domain $D(S) = D$: $Sy = y^*$. $S$ is continuous and $\|Sy\| \leqq \delta^{-1} \|y\|$, $y \in D(S)$, because 
\[
\delta \left\| S y \right\| ^{2} \leqq B (S y, S y) = (S y, y) \leqq \left\| S y \right\| \cdot \left\| y \right\|. \]

Moreover, $D = D(S)$ is a closed linear subspace of $X$. Proof: if $y_{n} \in D(S)$ and $\slim\limits_{n \to \infty} y_{n} = y_{\infty}$, then, by the continuity of $S$ proved above, $\{Sy_{n}\}$ is a Cauchy sequence and so has a limit $z = \slim Sy_{n}$. By the continuity of the scalar product, we have $\lim\limits_{n \to \infty} (x, y_{n}) = (x, y_{\infty})$. We have also, by \eqref{eq:III.7.2}, $\lim\limits_{n\to \infty}B(x,Sy_n) = B(x,z)$. Thus, by $(x,y_{n}) = B(x,Sy_{n})$, we must have $(x,y_{\infty}) = B(x,z)$ which proves that $y_{\infty}\in D$ and $Sy_{\infty} = z$. 

Therefore the first part of the Theorem, that is, the existence of the operator $S$ is proved if we can show that $D(S) = X$. Suppose $D(S) \neq X$. Then there exists a $w_0 \in X$ such that $w_0 \neq 0$ and $w_0 \in D(S)^\perp$. Consider the linear functional $F(z) = B(z, w_0)$ defined on $X$. $F(z)$ is continuous, since $|F(z)| = |B(z, w_0)| \leqq \gamma \|z\| \cdot \|w_0\|$. Thus, by F. \textsc{Riesz}' representation theorem, there exists a $w_0' \in X$ such that $B(z, w_0) = F(z) = (z, w_0')$ for all $z \in X$. This proves that $w_0' \in D(S)$ and $Sw_0' = w_0$. But, by $\delta \|w_0\|^2 \leqq B(w_0, w_0) = (w_0, w_0') = 0$, we obtain $w_0 = 0$ which is a contradiction. 

The inverse $S^{-1}$ exists. For, $Sy = 0$ implies $(x, y) = B(x, Sy) = 0$ for all $x \in X$ and so $y = 0$. As above, we show that, for every $y \in X$ there exists a $y'$ such that $(z, y') = B(z, y)$ for all $z \in X$. Hence $y = Sy'$ and so $S^{-1}$ is an everywhere defined operator, and, by $|(z, S^{-1}y)| = |B(z, y)| \leqq \gamma \|z\| \cdot \|y\|$, we see that $\|S^{-1}\| \leqq \gamma$. 
\end{proof}

Concrete applications of the Lax-Milgram theorem will be given in later chapters. In the following four sections, we shall give some examples of the direct application of F. \textsc{Riesz}' representation theorem. 


\yossection{8}{A Proof of the Lebesgue-Nikodym Theorem}
\label{sec:III.8}

This theorem reads as follows. 

\begin{theoremn}[\textsc{Lebesgue}-\textsc{Nikodym}] Let $(S, \mathfrak{B}, m)$ be a measure space, and $\nu(B)$ be a $\sigma$-finite, $\sigma$-additive and non-negative measure defined on $\mathfrak{B}$. If $\nu$ is $m$-absolutely continuous, then there exists a non-negative, $m$-measurable function $p(s)$ such that
\begin{equation}\label{eq:III.8.1}
\nu (B) = \int_{B} p (s) m (\dd s) \quad \text{ for all } \quad B \in \mathfrak {B} \text{ with } \nu (B) < \infty .
\end{equation}

Moreover, the ``density'' $p(s)$ of $v(B)$ (with respect to $m(B)$) is uniquely determined in the sense that any two of them are equal $m$-a.e. 
\end{theoremn}

Proof (due to J. von \textsc{Neumann} \cite{neumann229}). It is easy to see that $\varrho(B) = m(B) + \nu(B)$ is a $\sigma$-finite, $\sigma$-additive and non-negative measure defined on $\mathfrak{B}$. Let $\{B_n\}$ be a sequence of sets $\in \mathfrak{B}$ such that $S = \bigcup\limits_{n=1}^{\infty} B_n$, $B_n \subseteq B_{n+1}$ and $\varrho(B_n) < \infty$ for $n = 1, 2, \ldots$. If we can prove the theorem for every $B \subseteq B_n$ (for fixed $n$) and obtain the density $\phi_n(s)$, then the Theorem is true. For, we have only to take $\phi(s)$ as follows: 
\begin{equation}p (s) = p _{1} (s) \text{ for } s \in B _{1}, \text{ and } p (s) = p _{n+1} (s) \text{ for } s \in B _{n+1} - B _{n} (n = 1, 2, \dots). \notag\end{equation}

Therefore we may assume that $\varrho(S) < \infty$ without losing the generality. Consider the Hilbert space $L^2(S, \mathfrak{B}, \varrho)$. Then 
\[
f (x) = \int_{S} x (s) \nu (\dd s), x \in L ^{2} (S, \mathfrak {B}, \varrho),
\]
gives a bounded linear functional on $L^{2}(S, \mathfrak{B}, \varrho)$, because 
\[
\begin{aligned}
&| f (x) |  \leqq \int_{S} | x (s) | \nu (\dd s) \leqq \left(\int_{S} | x (s) | ^{2} \nu (\dd s)\right) ^{1/2} \left(\int_{S} 1 \cdot \nu (\dd s)\right) ^{1/2} \leqq \| x \| _{\varrho} \cdot \nu (S) ^{1/2},
\end{aligned}
\]
where $\| x\| _{\varrho} = \left(\int \limits_{S}|x(s)|^2\varrho (\dd s)\right)^{1/2}$. Thus, by F. \textsc{Riesz}' representation theorem, there exists a uniquely determined $y\in L^{2}(S,\mathfrak{B},\varrho)$ such that 
\[
\int_{S} x (s) \nu (\dd s) = \int_{S} x (s) \overline {{y (s)}} \varrho (\dd s) = \int_{S} x (s) \overline {{y (s)}} m (\dd s) + \int_{S} x (s) \overline {{y (s)}} \nu (\dd s)
\]
holds for all $x \in L^{2}(S, \mathfrak{B}, \varrho)$. Taking $x$ as non-negative functions and considering the real part of both sides, we may assume that $y(s)$ is a real-valued function. Hence 
\begin{equation}\label{eq:III.8.2}
\int_{S} x (s) (1 - y (s)) v (\dd s) = \int_{S} x (s) y (s) m (\dd s) \text{ if } x (s) \in L ^{2} (S, \mathfrak {B}, \varrho)
\end{equation}
is non-negative. 

We can prove $0 \leqq y(s) < 1$ $\varrho$-a.e. To this purpose, set $E_{1} = \{s; y(s) < 0\}$ and $E_{2} = \{s; y(s) \geqq 1\}$. If we take the defining function $C_{E_1}(s)$ of $E_{1}$ for $x(s)$ in \eqref{eq:III.8.2}, then the left hand side is $\geqq 0$ and hence $\int_{E_1} y(s) m(\dd s) \geqq 0$. 

Thus we must have $m(E_1) = 0$, and so, by the $m$-absolute continuity of $\nu$, $\nu(E_1) = 0$, $\varrho(E_1) = 0$. We may also prove $\varrho(E_2) = 0$, by taking the defining function $C_{E_1}(s)$ for $x(s)$ in \eqref{eq:III.8.2}. Therefore $0 \leqq y(s) < 1$ $\varrho$-a.e. on $S$. 

Let $x(s)$ be $\mathfrak{B}$-measurable and $\geqq 0 \varrho$-a.e. Then, by $\varrho(S) < \infty$, the ``truncated functions'' $x_{n}(s) = \min(x(s), n)$ belong to $L^{2}(S, \mathfrak{B}, \varrho) (n = 1, 2, \ldots)$, and hence 
\begin{equation}\label{eq:III.8.3}
\int_{S} x _{n} (s) (1 - y (s)) \nu (\dd s) = \int_{S} x _{n} (s) y (s) m (\dd s) \quad (n = 1, 2, \dots).
\end{equation}

Since the integrals increase monotonely as $n$ increases, we have 
\begin{equation}\label{eq:III.8.4}
\lim\limits _{n \rightarrow \infty} \int_{S} x _{n} (s) (1 - y (s)) v (\dd s) = \lim\limits _{n \rightarrow \infty} \int_{S} x _{n} (s) y (s) m (\dd s) = L \leqq \infty .
\end{equation}

Since the integrands are $\geqq 0$ $\varrho$-a.e., we have, by the Lebesgue-Fatou Lemma, 
\begin{align}
L& \geqq \int_{S} \varliminf\limits_{n \rightarrow \infty} (x _{n} (s) (1 - y (s)) v (\dd s) = \int_{S} x (s) (1 - y (s)) v (\dd s), \label{eq:III.8.5} \\
L& \geqq \int_ {\mathcal {S}} \lim\limits _{n \rightarrow \infty} (x _{n} (s) y (s)) m (\dd s) = \int_ {\mathcal {S}} x (s) y m (\dd s), \notag
\end{align}
under the convention that if $x(s) (1 - y(s))$ is not $\nu$-integrable, the right hand side is equal to $\infty$; and the same convention for $x(s) y(s)$. If $x(s) y(s)$ is $m$-integrable, then, by the Lebesgue-Fatou Lemma, 
\begin{equation}\label{eq:III.8.6}
L \leqq \int_{S} \overline {{\lim\limits _{n \rightarrow \infty}}} (x _{n} (s) y (s)) m (\dd s) = \int_{S} x (s) y (s) m (\dd s).
\end{equation}

This formula is true even if $x(s)y(s)$ is not $m$-integrable, under the convention that then $L=\infty$. Under a similar convention, we have 
\begin{equation}\label{eq:III.8.7}
L \leqq \int_{S} x (s) (1 - y (s)) v (\dd s).
\end{equation}

Therefore, we have 
\begin{equation}\label{eq:III.8.8}
\int_{S} x (s) (1 - y (s)) v (\dd s) = \int_{S} x (s) y (s) m (\dd s) \text{ for every } x (s) \text{ which }
\end{equation}
is $\mathfrak{B}$-measurable and $\geqq 0$ $\varrho$-a.e., 
under the convention that, if either side of the equality is $=\infty$ then the other side is also $=\infty$. 

Now we put 
\[
x (s) (1 - y (s)) = z (s), y (s) (1 - y (s)) ^{-1} = p (s).
\]

Then, under the same convention as in \eqref{eq:III.8.8}, we have 
\begin{equation}\label{eq:III.8.9}
\int_{S} z (s) \nu (\dd s) = \int_{S} z (s) p (s) m (\dd s) \text{ if } z (s) \text{ is } \mathfrak {B} \text{-measurable }
\end{equation}
and $\geqq 0$ $\varrho$-a.e. 

If we take the defining function $C_B(s)$ of $B \in \mathfrak{B}$ for $z(s)$, we obtain 
\[
\nu (B) = \int_{B} p (s) m (\dd s) \text{ for all } B \in \mathfrak {B}.
\]

The last part of the Theorem is clear in view of definition \eqref{eq:III.8.1}. 

Reference. For a straightforward proof of the Lebesgue-Nikodym theorem based upon Hahn's decomposition (Theorem~\ref{thm:I.3.Theorem.3} in Chapter~\ref{chap:I}, \ref{sec:I.3}) see K. \textsc{Yosida} \cite{yosida345}. This proof is reproduced in \textsc{Halmos} \cite{halmos92}, p. $128$. See also \textsc{Saks} \cite{saks280} and \textsc{Dunford}-\textsc{Schwartz} \cite{dunford44}. 


\yossection{9}{The Aronszajn-Bergman Reproducing Kernel}
\label{sec:III.9}

Let $A$ be an abstract set, and let a system $X$ of complex-valued functions defined on $A$ constitute a Hilbert space by the scalar product 
\begin{equation}\label{eq:III.9.1}
(f, g) = (f (a), g (a)) _{a}.
\end{equation}

A complex-valued function $K(a, b)$ defined on $A \times A$ is called a reproducing kernel of $X$ if it satisfies the condition: 
\begin{align}
\text{ For any fixed } b, K (a, b)& \in X \text{ as a function of } a, \label{eq:III.9.2} \\
f (b)& = \left(f (a), K (a, b)\right) _{a} \text{ and hence } \overline {{f (b)}} = \left(K (a, b), f (a)\right) _{a}. \label{eq:III.9.3}
\end{align}

As for the existence of reproducing kernels, we have 

\begin{theorem}[N. \textsc{Aronszajn} {\cite{aronszajn7}}, S. \textsc{Bergman} {\cite{bergman15}}]\label{thm:III.9.Theorem.1} $X$ has a reproducing kernel $K$ iff there exists, for any $y_{0} \in A$, a positive constant $C_{y_{0}}$, depending upon $y_{0}$, such that 
\begin{equation}\label{eq:III.9.4}
\left| f \left(y _{0}\right) \right| \leqq C _ {y _{0}} \left\| f \right\| \quad \text{ for all } \quad f \in X.
\end{equation}
\end{theorem}

\begin{proof}
 The ``only if'' part is proved by applying Schwarz' inequality to $f(y_{0}) = (f(x), K(x, y_{0}))_{x}$: 
\begin{equation}\label{eq:III.9.5}
\left| f \left(y _{0}\right) \right| \leqq \left\| f \right\| \cdot \left(K \left(x, y _{0}\right), K \left(x, y _{0}\right)\right) _{x} ^{1/2} = \left\| f \right\| K \left(y _{0}, y _{0}\right) ^{1/2}. \end{equation}

The ``if'' part is proved by applying F. \textsc{Riesz}' representation theorem to the linear functional $F_{y_{0}}(f) = f(y_{0})$ of $f \in X$. Thus, there exists a uniquely determined vector $g_{y_{0}}(x)$ of $X$ such that, for every $f \in X$, 
\[
f (y _{0}) = F _ {y _{0}} (f) = (f (x), g _ {y _{0}} (x)) _{x},
\]
and so $g_{y_0}(x) = K(x, y_0)$ is a reproducing kernel of $X$. The proof shows that the reproducing kernel is uniquely determined. 
\end{proof}

\begin{corollaryu} We have 
\begin{equation}\label{eq:III.9.6}
\sup\limits _{\| f \| \leqq 1} | f (y _{0}) | = K (y _{0}, y _{0}) ^{1/2},
\end{equation}
the supremum being attained by 
\begin{equation}\label{eq:III.9.7}
f _{0} (x) = \varrho K (x, y _{0}) / K (y _{0}, y _{0}) ^{1/2}, | \varrho | = 1.
\end{equation}
\end{corollaryu}

\begin{proof}
 The equality in the Schwarz' inequality \eqref{eq:III.9.5} holds iff $f(x)$ and $K(x, y_0)$ are linearly dependent. From the two conditions $f(x) = \alpha K(x, y_0)$ and $\|f\| = 1$, we obtain 
\[
1 = | \alpha | (K (x, y _{0}), K (x, y _{0})) _{x} ^{1/2} = | \alpha | K (y _{0}, y _{0}) ^{1/2}, \text{ that is, } | \alpha | = K (y _{0}, y _{0}) ^{-1/2}. \]

Hence the equality sign in \eqref{eq:III.9.5} is attained by $f_0(x)$. 
\end{proof}

\begin{examplen} Consider the Hilbert space $A^2(G)$. For any $f \in A^2(G)$ and $z \in G$, we have (see \eqref{eq:I.9.4} in Chapter~\ref{chap:I}, \ref{sec:I.9})
\[
\left| f \left(z _{0}\right) \right| ^{2} \leqq (\pi r ^{2}) ^{-1} \int_ {\left| z - z _{0} \right| \leqq r} | f (z) | ^{2} \dd x \dd y (z = x + i y).
\]

Thus $A^2(G)$ has the reproducing kernel which will be denoted by $K_G(z, z')$. This $K_G(z, z')$ is called Bergman's kernel of the domain $G$ of the complex plane. The following theorem of Bergman illustrates the meaning of $K_G(z, z')$ in the theory of conformal mapping. 

\end{examplen}
\begin{theorem}\label{thm:III.9.Theorem.2} Let $G$ be a simply connected bounded open domain of the complex plane, and $z_{0}$ be any point of $G$. By Riemann's theorem, there exists a uniquely determined regular function $w = f(z; z_{0})$ of $z$ which gives a one-to-one conformal map of the domain $G$ onto the sphere $|w| < \varrho_{G}$ of the complex $w$-plane in such a way that 
\[
f _{0} (z _{0}; z _{0}) = 0, (\dd f _{0} (z; z _{0}) / \dd z) _ {z = z _{0}} = 1.
\]
\end{theorem}

Bergman's kernel $K_{G}(z;z_{0})$ is connected with $f_{0}(z;z_{0})$ by 
\begin{equation}\label{eq:III.9.8}
f _{0} (z; z _{0}) = K _{G} \left(z _{0}; z _{0}\right) ^{-1} \int_ {z _{0}} ^{z} K _{G} (t; z _{0}) \dd t,
\end{equation}
where the integral is taken along any rectifiable curve lying in $G$ and connecting $z_{0}$ with $z$. 

$A_{1}^{2}(G) = \{f(z); f(z)\text{ is holomorphic in } G, f'(z) \in A^{2}(G), f(z_{0}) = 0\text{ and } f'(z_{0}) = 1\}$, and consider, for any $f \in A_{1}^{2}(G)$, the number 
\begin{equation}\label{eq:III.9.9}
\left\| f ^{\prime} \right\| ^{2} = \int_{G} \left| f ^{\prime} (z) \right| ^{2} \dd x \dd y, z = x + i y.
\end{equation}

If we denote by $z = \varphi(w)$ the inverse function of $w = f_0(z; z_0)$, then, for any $f \in A_1^2(G)$, 
\[
\left\| f ^{\prime} \right\| ^{2} = \iint_ {| w | < \varrho_{G}} \left| f ^{\prime} (\varphi (w)) \right| ^{2} \left| \varphi^{\prime} (w) \right| ^{2} \dd u \dd v, w = u + i v.
\]

For, by the Cauchy-Riemann partial differential equations 
\[
x _{u} = y _{v}, x _{v} = - y _{u},
\]
we have 
\[
\begin{aligned}
&\dd x \dd y  = \frac {\partial (x , y)}{\partial (u , v)} \dd u \dd v = (x _{u} y _{v} - y _{u} x _{v}) \dd u \dd v = (x _{u} ^{2} + y _{u} ^{2}) \dd u \dd v = | \varphi^{\prime} (w) | ^{2} \dd u \dd v.
\end{aligned}
\]

Let $f \in A_1^2(G)$, and let $F(w) = f(\varphi(w))$ be expanded into power series: 
\[
F (w) = f (\varphi (w)) = w + \sum\limits_{n=2} ^{\infty} c _{n} w ^{n} \text{ for } | w | < \varrho_{G}.
\]

Then $F'(w) = f'(\varphi(w)) \varphi'(w) = 1 + \sum\limits_{n=2}^{\infty} n c_n w^{n-1}$ and so 
\[
\begin{aligned}
\left\| f ^{\prime} \right\| ^{2} &= \iint_ {| w | < \varrho_{G}} \left| 1 + \sum\limits_{n=2} ^{\infty} n c _{n} w ^{n-1} \right| ^{2} \dd u \dd v
\\
&= \int_{0} ^ {\varrho_{G}} \dd r \left\{\int_{0} ^{2 \pi} \left(r + \sum\limits_{n=2} ^{\infty} n ^{2} \left| c _{n} \right| ^{2} r ^{2n-1}\right) \dd \theta \right\} = \pi \varrho_{G} ^{2} + \sum\limits_{n=2} ^{\infty} \pi n \left| c _{n} \right| ^{2} \varrho_{G} ^{2n}.
\end{aligned}
\]

Therefore minimum $\left\|f^{\prime}\right\|=\sqrt{\pi}\varrho_{G}$, and this minimum is attained iff $F(w)=f(\varphi(w))=w$, that is, iff $f(z)=f_{0}(z;z_{0})$. 

We set, for any $f \in A_1^2(G)$, $g(z) = f(z) / \|f'\|$. Then $\|g'\| = 1$. If we put $\tilde{A}^2(G) = \{g(z); g(z) \text{ is holomorphic in } G, g(z_0) = 0, g'(z_0) > 0 \text{ and } \|g'\| = 1\}$, 
then the above remark shows that 
\[
\underset {g \in \widetilde {A} ^{2} (G)} {\text{ maximum }} g ^{\prime} \left(z _{0}\right) = 1 / \left\| f _{0} ^{\prime} \right\| = (\sqrt {\pi} \varrho_{G}) ^{-1},
\]
and this maximum is attained iff $g(z)$ is equal to 
\[
g _{0} (z) = f _{0} \left(z; z _{0}\right) / \left\| f _{0} ^{\prime} \right\| = f _{0} \left(z; z _{0}\right) / \sqrt {\pi} \varrho_{G}.
\]

Hence, by \eqref{eq:III.9.7}, we obtain 
\[
g _{0} ^{\prime} (z) = (\sqrt {\pi} \varrho_{G}) ^{-1} \frac {\dd f _{0} (z ; z _{0})}{\dd z} = \lambda K _{G} (z; z _{0}) / K _{G} (z _{0}; z _{0}) ^{1/2}, | \lambda | = 1.
\]

Hence, by putting $z = z_0$, we have 
\[
\left(\lambda \sqrt {\pi} \varrho_{G}\right) ^{-1} = K _{G} (z _{0}; z _{0}) / K _{G} (z _{0}; z _{0}) ^{1/2} = K _{G} (z _{0}; z _{0}) ^{1/2},
\]
and so we have proved the formula 
\[
\frac {\dd f _{0} (z ; z _{0})}{\dd z} = K _{G} (z; z _{0}) / K _{G} (z _{0}; z _{0}).
\]


\yossection{10}{The Negative Norm of P. \textsc{Lax}}
\label{sec:III.10}

Let $H_{0}^{s}(\Omega)$ be the completion of the pre-Hilbert space $C_{0}^{\infty}(\Omega)$ endowed with the scalar product $(\varphi, \psi)_{s}$ and the norm $\| \varphi \|_{s}$: 
\begin{equation}\label{eq:III.10.1}
(\varphi , \psi) _{s} = \sum\limits_{| j | \leqq s} \int_{\Omega} D ^{j} \varphi (x) \overline {{D ^{j} \psi (x)}} \dd x, \| \varphi \| _{s} = (\varphi , \varphi) _{s} ^{1/2}.
\end{equation}

Any element $b \in H_0^0(\Omega) = L^2(\Omega)$ defines a continuous linear functional $f_b$ on $H_0^s(\Omega)$ by 
\begin{equation}\label{eq:III.10.2}
f _{b} (w) = (w, b) _ {\mathbf {0}}, w \in H _ {\mathbf {0}} ^{s} (\Omega).
\end{equation}

For, by Schwarz' inequality, we have 
\[
\left| (w, b) _{0} \right| \leqq \left\| w \right\| _{0} \cdot \left\| b \right\| _{0} \leqq \left\| w \right\| _{s} \cdot \left\| b \right\| _{0}.
\]

Therefore, if we define the negative norm of $b \in H_0^0(\Omega) = L^2(\Omega)$ 
\begin{equation}\label{eq:III.10.3}
\left\| b \right\| _{-s} = \sup\limits _ {w \in H _{0} ^{\varepsilon} (\Omega), \| w \| _{s} \leqq 1} \left| f _{b} (w) \right| = \sup\limits _ {w \in H _{0} ^{\varepsilon} (\Omega), \| w \| _{s} \leqq 1} \left| (w, b) _{0} \right|,
\end{equation}
then we have 
\begin{equation}\label{eq:III.10.4}
\left\| b \right\| _{-s} \leqq \left\| b \right\| _{0},
\end{equation}
and so, by $\|b\|_{-s} \geqq |(w/\|w\|_s, b)_0|$, 
\begin{equation}\label{eq:III.10.5}
\left| (w, b) _{0} \right| \leqq \left\| w \right\| _{s} \cdot \left\| b \right\| _{-s}.
\end{equation}

Hence we may write 
\begin{equation}\tag{3'}\label{eq:III.10.3'}
\left\| b \right\| _{-s} = \left\| f _{b} \right\| _{-s} = \sup\limits _ {\left\| w \right\| _{s} \leqq 1} | (w, b) _{0} | \text{ for any } b \in H _{0} ^{0} (\Omega).
\end{equation}

We shall prove 

\begin{theorem}[P. \textsc{Lax} {\cite{lewy172}}]\label{thm:III.10.Theorem.1} The dual space $H_0^s(\Omega)'$ of the space $H_0^s(\Omega)$ may be identified with the completion of the space $H_0^0(\Omega) = L^2(\Omega)$ with respect to the negative norm. 
\end{theorem}

For the proof we prepare 

\begin{propositionu} The totality $F$ of the continuous linear functionals on $H_0^s(\Omega)$ of the form $f_b$ is dense in the Hilbert space $H_0^s(\Omega)'$, which is the dual space of $H_0^s(\Omega)$. 
\end{propositionu}

\begin{proof}
 $F$ is total on $H_0^s(\Omega)$ in the sense that, for a fixed $w \in H_0^s(\Omega)$, the simultaneous vanishing $f_b(w) = 0$, $b \in H_0^0(\Omega)$, takes place only if $w = 0$. This is clear since any element $w \in H_0^s(\Omega)$ is also an element of $H_0^0(\Omega)$. 

Now if $F$ is not dense in the Hilbert space $H_0^s(\Omega)'$, then there exists an element $T \neq 0$ of the second dual space $H_0^s(\Omega)'' = (H_0^s(\Omega)')'$ such that $T(f_b) = 0$ for all $f_b \in F$. By the reflexivity of the Hilbert space $H_0^s(\Omega)$, there exists an element $t \in H_0^s(\Omega)$ such that $T(f) = f(t)$ for all $f \in H_0^s(\Omega)'$. Thus $T(f_b) = f_b(t) = 0$ for all $b \in H_0^0(\Omega)$. This implies, by the totality of $F$ proved above, that $t = 0$, contrary to $T \neq 0$. 
\end{proof}
\begin{corollaryu} We have, dually to \eqref{eq:III.10.3}, 
\begin{equation}\label{eq:III.10.6}
\left\| w \right\| _{s} = \sup\limits _ {b \in H _{0} ^{0} (\Omega), \| b \| _{-s} \leqq 1} \left| (w, b) _{0} \right| \text{ for any } w \in H _{0} ^{s} (\Omega).
\end{equation}
\end{corollaryu}

\begin{proof}
 Clear from Corollary~\ref{cor:III.6.Corollary.3} in Chapter~\ref{chap:III}, \ref{sec:III.6}, because $F = \{f_b; b \in H_0^0(\Omega)\}$ is dense in $H_0^s(\Omega)'$. 
\end{proof}

\begin{proof}[Proof of Theorem~\ref{thm:III.10.Theorem.1}.]
Clear from the facts that \begin{enumerate}[i)]
\item $F$ is dense in the dual space $H_{0}^{s}(\Omega)'$
\item $F$ is in one-to-one correspondence to the set $H_{0}^{0}(\Omega)=L^{2}(\Omega)$ preserving the negative norm, i.e.,
\[
F \ni f _{b} \leftrightarrow b \in H _{0} ^{0} (\Omega) \text{ and } \| f _{b} \| _{-s} = \| b \| _{-s}.
\]
\end{enumerate}


We shall denote by $H_0^{-s}(\Omega)$ the completion of $H_0^0 (\Omega)$ with respect to the negative norm $\|b\|_{-s}$. Thus 
\begin{equation}\label{eq:III.10.7}
H _{0} ^{s} (\Omega) ^{\prime} = H _{0} ^{-s} (\Omega).
\end{equation}

For any continuous linear functional $f$ on $H_0^s(\Omega)$, we shall denote by $\langle w, f \rangle$ the value of $f$ at $w \in H_0^s(\Omega)$. Thus, for any $b \in H_0^0(\Omega)$, we may write 
\begin{equation}\label{eq:III.10.8}
f _{b} (w) = (w, b) _{0} = \langle w, f _{b} \rangle = \langle w, b \rangle , w \in H _{0} ^{s} (\Omega),
\end{equation}
and have the generalized Schwarz' inequality 
\begin{equation}\label{eq:III.10.9}
\left| \langle w, b \rangle \right| \leqq \left\| w \right\| _{s} \left\| b \right\| _{-s},
\end{equation}
which is precisely \eqref{eq:III.10.5}. 

Now we can prove 
\end{proof}

\begin{theorem}[P. \textsc{Lax} {\cite{lewy172}}]\label{thm:III.10.Theorem.2} Any continuous linear functional $g(b)$ on $H_0^{-s}(\Omega)$ can be represented, by a fixed element $w \in H_0^s(\Omega)$, as 
\begin{equation}\label{eq:III.10.10}
g (b) = g _{w} (b) = \overline {{\langle w , b \rangle}}.
\end{equation}
\end{theorem}

We have, in particular, 
\begin{equation}\label{eq:III.10.11}
H _{0} ^{s} (\Omega) ^{\prime} = H _{0} ^{-s} (\Omega), H _{0} ^{-s} (\Omega) ^{\prime} = H _{0} ^{s} (\Omega).
\end{equation}



\begin{proof}
 If $b \in H_0^0(\Omega)$, then $\langle w, b \rangle = f_b(w) = (w, b)_0$. Since $F = \{f_b; b \in H_0^0(\Omega)\}$ is dense in the Hilbert space $H_0^s(\Omega)'$, we know, by \eqref{eq:III.10.9}, that $\overline{\langle w, b \rangle} = (b, w)_0$ with a fixed $w \in H_0^s(\Omega)$, defines a linear functional $g_w$ continuous on a dense subset $F$ of $H_0^s(\Omega)'$. The norm of this functional $g_w$ on $F$ will be denoted by $\|g_w\|_s$. Then, by \eqref{eq:III.10.6}, 
\begin{equation}\label{eq:III.10.12}
\left\| g _{w} \right\| _{s} = \sup\limits _{\left\| b \right\| - s \leqq 1} \left| (b, w) _{0} \right| = \sup\limits _{\left\| b \right\| - s \leqq 1} \left| (w, b) _{0} \right| = \left\| w \right\| _{s}. \end{equation}

We may thus extend, by continuity, the functional $g_{w}$ on $F$ to a continuous linear functional on the completion of $F$ (with respect to the negative norm), that is, $g_{w}$ can be extended to a continuous linear functional on $H_{0}^{s}(\Omega)' = H_{0}^{-s}(\Omega)$; we denote this extension by the same letter $g_{w}$. We have thus 
\begin{equation}\label{eq:III.10.13}
\left\| g _{w} \right\| = \sup\limits _{\left\| b \right\| - s \leqq 1} \left| g _{w} (b) \right| = \left\| w \right\| _{s}.
\end{equation}

Hence, in view of the completeness of the space $H_0^s(\Omega)$, the totality $G$ of the continuous linear functionals $g_w$ on $H_0^{-s}(\Omega)$ may be considered as a closed linear subspace of $H_0^{-s}(\Omega)'$ by the correspondence $g_w \leftrightarrow w$. If this closed subspace $G$ were not dense in $H_0^{-s}(\Omega)'$, then there exists a continuous linear functional $f \neq 0$ on $H_0^{-s}(\Omega)'$ such that $f(g_w) = 0$ for all $g_w \in G$. But, since the Hilbert space $H_0^{-s}(\Omega)$ is reflexive, such a functional $f$ is given by $f(g_w) = g_w(f_0)$, $f_0 \in H_0^{-s}(\Omega)$, and so by \eqref{eq:III.10.13} $f_0$ must be equal to $0$, contrary to the fact $f \neq 0$. Therefore we have proved $H_0^{-s}(\Omega)' = H_0^s(\Omega)$. 
\end{proof}

\begin{remarku} The notion of the negative norm was introduced by P. \textsc{Lax} with the view of applying it to the genuine differentiability of distribution solutions of linear partial differential equations. We shall discuss such differentiability problems in later chapters. It is to be noted that the notion of the negative norm is also introduced naturally through the Fourier transform. This was done by J. \textsc{Leray} \cite{leray175} earlier than \textsc{Lax}. We shall explain the point in a later chapter on the Fourier transform. 
\end{remarku}


\yossection{11}{Local Structures of Generalized Functions}
\label{sec:III.11}

A generalized function is locally the distributional derivative of a function. More precisely, we can prove 

\begin{theoremn}[L. \textsc{Schwartz} {\cite{schwartz294}}] Let $T$ be a generalized function in $\Omega \subseteq \mathbb{R}^{n}$. Then, for any compact subset $K$ of $\Omega$, there exist a positive integer $m_0 = m_0(T, K)$ and a function $f(x) = f(x; T, K, m_0) \in L^2(K)$ such that
\begin{equation}\label{eq:III.11.1}
T (\varphi) = \int_{K} f (x) \frac {\partial^ {n m _{0}} \varphi (x)}{\partial x _{1} ^ {m _{0}} \partial x _{2} ^ {m _{0}} \cdots \partial x _{n} ^ {m _{0}}} \dd x \text{ whenever } \varphi \in \mathfrak {D} _{K} (\Omega)
\end{equation}

\end{theoremn}
\begin{proof}
 By the Corollary in Chapter~\ref{chap:I}, \ref{sec:I.8}, there exist a positive constant $C$ and a positive integer $m$ such that 
\begin{equation}\label{eq:III.11.2}
\left| T (\varphi) \right| \leqq C \sup\limits _{| j | \leqq m, x \in K} \left| D ^{j} \varphi (x) \right| \text{ whenever } \varphi \in \mathfrak {D} _{K} (\Omega). \end{equation}

Thus there exists a positive number $\delta$ such that 
\begin{equation}\label{eq:III.11.3}
\phi_{m} (\varphi) = \sup\limits _{| j | \leqq m, x \in K} | D ^{j} \varphi (x) | \leqq \delta , \varphi \in \mathfrak {D} _{K} (\Omega), \text{ implies } | T (\varphi) | \leqq 1.
\end{equation}

We introduce the notation 
\begin{equation}\label{eq:III.11.4}
\frac {\partial^{s}}{\partial x ^{s}} = \frac {\partial^{sn}}{\partial x _{1} ^{s} \partial x _{2} ^{s} \cdots \partial x _{n} ^{s}},
\end{equation}
and prove that there exists a positive constant $\varepsilon$ such that, for $m_0 = m + 1$, 
\begin{equation}\label{eq:III.11.5}
\int_{\Omega} \left| \partial^ {m _{0}} \varphi (x) / \partial x ^ {m _{0}} \right| ^{2} \dd x \leqq \varepsilon , \varphi \in \mathfrak {D} _{K} (\Omega), \text{ implies } p _{m} (\varphi) \leqq \delta .
\end{equation}

This is proved by repeated application of the following inequality 
\[
\begin{aligned}
\left| \psi (x) \right|& \leqq \underset {K \cap (- \infty , x _{i})} {\int} \left| \partial \psi (x _{1}, \dots , x _{i-1}, y, x _{i+1}, \dots , x _{n}) / \partial y \right| \dd y
\\
&\leqq \left(\underset {K \cap (- \infty , x _{i})} {\int} \dd y\right) ^{1/2} \left(\underset {K \cap (- \infty , x _{i})} {\int} \left| \partial \psi (x _{1}, \dots , x _{i-1}, y, x _{i+1}, \dots , x _{n}) / \partial y \right| ^{2} \dd y\right) ^{1/2}
\\
&= t ^{1/2} \left(\underset {K \cap (- \infty , x _{i})} {\int} \left| \partial \psi (x _{1}, \dots , x _{i-1}, y, x _{i+1}, \dots , x _{n}) / \partial y \right| ^{2} \dd y\right) ^{1/2},
\end{aligned}
\]
where $t$ is the diameter of $K$, i.e., the maximum distance between two points of the compact set $K$. 

Consider the mapping $\varphi(x) \to \psi(x) = \partial^{m_0} \varphi(x)/\partial x^{m_0}$ defined on $\mathfrak{D}_K(\Omega)$ into $\mathfrak{D}_K(\Omega)$. As may be seen by integration, $\psi(x) = 0$ implies $\varphi(x) = 0$. Hence the above mapping is one-to-one. Thus $T(\varphi)$, $\varphi \in \mathfrak{D}_K(\Omega)$, defines a linear functional $S(\psi)$, $\psi(x) = \partial^{m_0} \varphi(x)/\partial x^{m_0}$, by $S(\psi) = T(\varphi)$. By \eqref{eq:III.11.3} and \eqref{eq:III.11.5}, $S$ is a continuous linear functional on the pre-Hilbert space $X$ consisting of such $\psi$'s and topologized by the norm $\|\psi\| = \left( \int_K |\psi(x)|^2 \dd x \right)^{1/2}$. 

Thus there exists, by F. \textsc{Riesz}' representation theorem, a uniquely determined function $f(x)$ from the completion of $X$ such that 
\[
T (\varphi) = S (\psi) = \int_{K} \left(\partial^ {m _{0}} \varphi (x) / \partial x ^ {m _{0}}\right) \cdot f (x) \dd x \text{ for all } \varphi \in \mathfrak {D} _{K} (\Omega).
\]

Actually, the completion of $X$ is contained in $L^{2}(K)$ as a closed linear subspace, and so the Theorem is proved. 
\end{proof}

\yossub{References for Chapter III}

For general account of Hilbert spaces, see N. I. \textsc{Achieser}-I. M. \textsc{Glazman} \cite{achieser1}, N. \textsc{Dunford}-J. \textsc{Schwartz} \cite{dunford45}, B. Sz. \textsc{Nagy} \cite{nagy217}, F. \textsc{Riesz}-B. Sz. \textsc{Nagy} \cite{nagy219} and M. H. \textsc{Stone} \cite{stone311}. 
